\documentclass[10pt,a4paper]{amsart}
\usepackage[margin=19mm]{geometry}
\usepackage{amsmath,amssymb,mathtools}
\usepackage[scr=rsfs,cal=euler]{mathalfa}
\usepackage{xfrac}
\usepackage[unicode,hidelinks,bookmarks=false]{hyperref}
\newcommand{\ie}{i.e.\ }
\newcommand{\cf}{cf.\ }
\newcommand{\resp}{respectively\ }
\newcommand{\Subsection}{\subsection}
\providecommand{\nobreakdash}{\nobreak\hbox{-}}
\setlength{\emergencystretch}{2em}
\multlinegap0pt
\usepackage{mathtools}
\usepackage{amssymb}
\usepackage[shortlabels]{enumitem}
\usepackage{xcolor}
\newtheorem{prop}{Proposition}
\newtheorem{definition}{Definition}
\DeclareMathOperator{\Tors}{Tors}
\title{Toral Chern--Simons TQFT\\via Geometric Quantization in Real Polarization}
\author{Daniel Galviz}
\date{Source: arXiv:2604.01016v1 (CC BY 4.0)}
\begin{document}
\maketitle

\noindent\emph{Source and licence.} Toral Chern--Simons TQFT\\via Geometric Quantization in Real Polarization. Version arXiv:2604.01016v1 (CC BY 4.0), distributed by arXiv under the Creative Commons Attribution 4.0 International licence, http://creativecommons.org/licenses/by/4.0/, as recorded in the arXiv licence metadata for this exact version and shown on the abstract page https://arxiv.org/abs/2604.01016v1. Full licence notice: \texttt{licenses/arxiv-2604.01016-CC-BY-4.0.txt}.\par\smallskip

\noindent\textit{Editorial scope.} Counted unit: the article's prop prop:Componentwise-toral-torsion-half-density (source0.tex lines 1784--1813). The article's complete original proof of it is delivered with it (source0.tex lines 1815--1935, one \texttt{proof} environment of 121 lines). No mathematics is written, repaired or re-derived: the statement and the proof are the article's own lines, in the article's own notation, and no retained line is substituted or re-flowed.  Retained uncounted as the source-local context the proof needs: lines 1592--1602; lines 1713--1747.  Nothing was omitted from any retained span, so the package carries no omission record.  No retained line contains a citation, so the wrapper carries no bibliography and no bibliography entry.  No retained line contains a figure environment, an image-inclusion command or drawing code, so no image is delivered and the package declares no asset dependency.  Editorial formatting only, in addition: an AMS \texttt{amsart} preamble that restates the article's own declarations and macros used by the retained lines, namely the environments \texttt{prop}, \texttt{definition} on separate counters, together with the standard packages those lines need.  The retained environments print in this document as Proposition 1, Definition 1 in order of appearance, whereas the article prints the same items with the numbering of its own full document.  The complete native source of the article as distributed in the arXiv e-print for this exact version is bundled unchanged as \texttt{sources/197628/source0.tex}.
\begin{prop}
\label{prop:toral-determinant-line}
For each $p \in \operatorname{Tors} H^2(X;\Lambda)$ there is a canonical isomorphism
\[
\bigl|\det(T_e^* \Lambda_{X,p})\bigr|^{1/2}
\cong
\bigl|\det(H^1(X;\mathfrak t)^*)\bigr|^{1/2}
\otimes
\bigl|\det(H^1(X,\partial X;\mathfrak t))\bigr|^{1/2}.
\]
\end{prop}

\begin{definition}
\label{def:Componentwise-half-density}
Fix $p\in \operatorname{Tors}H^2(X;\Lambda)$. Choose any nonzero element
\[
w\in \bigl|\det(H^1(X,\partial X;\mathfrak t)^*)\bigr|.
\]
Since $H^1(X,\partial X;\mathbb T)$ is a compact torus with tangent space
$H^1(X,\partial X;\mathfrak t)$ at the identity, the density $w$ determines a unique translation-invariant
density
\[
\overline{w}\in \Gamma\!\left(
H^1(X,\partial X;\mathbb T),\,
\bigl|\det T^* H^1(X,\partial X;\mathbb T)\bigr|
\right).
\]
Define
\begin{equation}
\label{eq:mu-tilde-definition}
\widetilde{\mu}_{X,p}
:=
\left(\int_{H^1(X,\partial X;\mathbb T)} \overline{w}\right)\,
T_X(\mathfrak t)^{1/2}\otimes w^{-1}.
\end{equation}

Using Proposition~\ref{prop:toral-determinant-line}, we regard $\widetilde{\mu}_{X,p}\in \bigl|\det(T_e^*\Lambda_{X,p})\bigr|^{1/2}.$ The corresponding translation-invariant half-density on $\Lambda_{X,p}$ is denoted
\[
\mu_{X,p}\in
\Gamma\!\left(
\Lambda_{X,p},\,
\bigl|\det(T^*\Lambda_{X,p})\bigr|^{1/2}
\right),
\]
and is obtained by transporting $\widetilde{\mu}_{X,p}$ from the identity tangent space to every point
of the leaf.
\end{definition}

\begin{prop}
\label{prop:Componentwise-toral-torsion-half-density}
For each $p\in\Tors H^2(X;\Lambda)$, the half-density $\mu_{X,p}$ of
Definition~\ref{def:Componentwise-half-density} satisfies:

\begin{enumerate}
\item[\textup{(i)}]
$\mu_{X,p}$ is natural under orientation-preserving diffeomorphisms of $X$;

\item[\textup{(ii)}]
for disjoint unions,
\[
\mu_{X_1\sqcup X_2,(p_1,p_2)}
=
\mu_{X_1,p_1}\boxtimes \mu_{X_2,p_2};
\]

\item[\textup{(iii)}]
if $t_\xi:\Lambda_{X,p}\to\Lambda_{X,p'}$ is translation by a vector between
parallel leaves, then
\[
\mu_{X,p'}=(t_\xi)_*\mu_{X,p};
\]

\item[\textup{(iv)}]
if $X$ is closed, the same construction yields the translation-invariant
density on $\mathcal M_{X,p}(\mathbb T)$ induced by the square root of
Ray--Singer/Reidemeister torsion.
\end{enumerate}
\end{prop}

\begin{proof}
\noindent
\emph{(i) Naturality.}
Let $f\colon X\to X'$ be an orientation-preserving diffeomorphism. Then $f$ induces canonical
isomorphisms
\[
f^*\colon H^1(X';\mathfrak t)\xrightarrow{\ \cong\ } H^1(X;\mathfrak t),
\qquad
f^*\colon H^1(X',\partial X';\mathfrak t)\xrightarrow{\ \cong\ } H^1(X,\partial X;\mathfrak t),
\]
and likewise on the corresponding determinant and density lines. It also identifies the torsion class
$p'\in \operatorname{Tors}H^2(X';\Lambda)$ with the corresponding class
$p=f^*p'\in \operatorname{Tors}H^2(X;\Lambda)$.

By the naturality of Reidemeister torsion under cellular subdivision and pullback by diffeomorphism, the
torsion density $T_{X'}(\mathfrak t)$ is carried to $T_X(\mathfrak t)$ under the induced determinant-line
isomorphism. Likewise, the canonical determinant-line identification of
Proposition~\ref{prop:toral-determinant-line} is functorial with respect to these induced maps. Therefore
the element
\[
\widetilde{\mu}_{X',p'}
\in \bigl|\det(T_e^*\Lambda_{X',p'})\bigr|^{1/2}
\]
of \eqref{eq:mu-tilde-definition} is transported to the corresponding element
\[
\widetilde{\mu}_{X,p}
\in \bigl|\det(T_e^*\Lambda_{X,p})\bigr|^{1/2}.
\]
Since $\mu_{X,p}$ and $\mu_{X',p'}$ are obtained by translation of these basepoint half-densities along
their respective leaves, it follows that $f$ identifies $\mu_{X',p'}$ with $\mu_{X,p}$. This proves
naturality.

\smallskip
\noindent
\emph{(ii) Disjoint unions.}
Let $X=X_1\sqcup X_2$, and let $p_i\in \operatorname{Tors}H^2(X_i;\Lambda)$ for $i=1,2$. Then
\[
H^1(X;\mathfrak t)\cong H^1(X_1;\mathfrak t)\oplus H^1(X_2;\mathfrak t),\,\,
H^1(X,\partial X;\mathfrak t)\cong
H^1(X_1,\partial X_1;\mathfrak t)\oplus H^1(X_2,\partial X_2;\mathfrak t),
\]
and similarly for the boundary leaves:
\[
\Lambda_{X,(p_1,p_2)}
=
\Lambda_{X_1,p_1}\times \Lambda_{X_2,p_2}.
\]
Under direct sum, determinant lines tensor-multiply canonically, and Reidemeister torsion is multiplicative.
Hence $T_X(\mathfrak t) = T_{X_1}(\mathfrak t)\otimes T_{X_2}(\mathfrak t)$
under the canonical identification of determinant lines. Likewise, if
\[
w_i\in \bigl|\det(H^1(X_i,\partial X_i;\mathfrak t)^*)\bigr|
\qquad (i=1,2)
\]
are nonzero densities, then
\[
w_1\otimes w_2
\in
\bigl|\det(H^1(X,\partial X;\mathfrak t)^*)\bigr|
\]
is the corresponding density for the disjoint union, and the induced Haar densities satisfy
\[
\overline{w_1\otimes w_2}
=
\overline{w}_1\boxtimes \overline{w}_2.
\]
Therefore
\[
\int_{H^1(X,\partial X;\mathbb T)} \overline{w_1\otimes w_2}
=
\left(\int_{H^1(X_1,\partial X_1;\mathbb T)} \overline{w}_1\right)
\left(\int_{H^1(X_2,\partial X_2;\mathbb T)} \overline{w}_2\right).
\]
Substituting these identities into \eqref{eq:mu-tilde-definition}, we obtain
\[
\widetilde{\mu}_{X,(p_1,p_2)}
=
\widetilde{\mu}_{X_1,p_1}\otimes \widetilde{\mu}_{X_2,p_2}.
\]
After extending these basepoint values to the corresponding translation-invariant half-densities on the leaves, this gives
\[
\mu_{X_1\sqcup X_2,(p_1,p_2)}
=
\mu_{X_1,p_1}\boxtimes \mu_{X_2,p_2},
\]
as claimed.

\smallskip
\noindent
\emph{(iii) Translation between parallel leaves.}
Fix two torsion classes $p,p'\in \operatorname{Tors}H^2(X;\Lambda)$. The leaves $\Lambda_{X,p}$ and
$\Lambda_{X,p'}$ are parallel translates of the same underlying Lagrangian torus, and
Proposition~\ref{prop:toral-determinant-line} identifies their tangent density lines with the same abstract
one-dimensional half-density line
\[
\bigl|\det(T_e^*\Lambda_{X,p})\bigr|^{1/2}
\cong
\bigl|\det(H^1(X;\mathfrak t)^*)\bigr|^{1/2}
\otimes
\bigl|\det(H^1(X,\partial X;\mathfrak t))\bigr|^{1/2}.
\]
In particular, the defining element $\widetilde{\mu}_{X,p}$ depends only on the cohomological data of
$X$ and not on the position of the leaf inside the ambient moduli space. Therefore the translation
$t_\xi\colon \Lambda_{X,p}\to \Lambda_{X,p'}$ carries the translation-invariant half-density determined by
$\widetilde{\mu}_{X,p}$ to the one determined by the same abstract tangent datum on $\Lambda_{X,p'}$.
Hence
\[
\mu_{X,p'}=(t_\xi)_*\mu_{X,p}.
\]

\smallskip
\noindent
\emph{(iv) Closed case.}
Assume now that $\partial X=\varnothing$. Then there is no boundary leaf; instead one works on the full
flat moduli torus $\mathcal M_{X,p}(\mathbb T).$ The same determinant-line formalism applies, but without relative terms. Reidemeister torsion now gives a
positive translation-invariant density on $\mathcal M_{X,p}(\mathbb T),$ or equivalently a positive element of the density line $\bigl|\det(H^1(X;\mathfrak t)^*)\bigr|.$
Taking its positive square root yields a translation-invariant half-density, and hence a translation-invariant
density on $\mathcal M_{X,p}(\mathbb T)$ induced by the square root of
Reidemeister\/--Ray\/--Singer torsion. This is precisely the closed analogue of the boundary-leaf
construction above.
\end{proof}

\end{document}
